%% Chapter 8 — Cross-Examination and Spot Check
\chapter{Cross-Examination and Spot Check}
\label{ch:spotcheck}

\section{Introduction}

Independent validation of the \BDE{} outputs is essential for establishing confidence in the reported delta values. This chapter presents a systematic cross-examination of the pipeline results against three independent data sources:

\begin{enumerate}
  \item Published professional ecological survey reports for sites within or adjacent to the study area;
  \item A University of Derby undergraduate thesis examining habitat condition in the Ashleyhay area;
  \item Natural England Priority Habitat Inventory (PHI) entries with associated condition data.
\end{enumerate}

The purpose of this cross-examination is to quantify the agreement between the \BDE{}'s remote sensing-derived classifications and condition assessments and independent ground-truth observations, identifying systematic biases or failure modes that may affect the reliability of the delta analysis.


\section{Validation Framework}

\subsection{Agreement Metrics}

For habitat classification validation, agreement is quantified using:

\begin{align}
\text{Overall Agreement (OA)} &= \frac{\text{Number of matching classifications}}{\text{Total comparison parcels}} \label{eq:oa-validation} \\[6pt]
\text{Cohen's Kappa} &= \frac{p_o - p_e}{1 - p_e} \label{eq:kappa}
\end{align}

\noindent where $p_o$ is the observed agreement and $p_e$ is the expected agreement by chance.

For condition assessment validation, agreement is measured as:

\begin{align}
\text{Exact Agreement} &= \frac{\text{Parcels with identical condition}}{\text{Total comparison parcels}} \label{eq:exact-cond} \\[6pt]
\text{Within-One Agreement} &= \frac{\text{Parcels within } \pm 1 \text{ category}}{\text{Total comparison parcels}} \label{eq:within-one}
\end{align}


\section{Professional Ecological Survey Data}

\subsection{Source Identification}

A search of publicly available ecological survey reports was conducted through:
\begin{itemize}[nosep]
  \item Amber Valley Borough Council planning register (ecology reports submitted with applications);
  \item Natural England site-specific monitoring reports for nearby SSSIs;
  \item Derbyshire Wildlife Trust survey data (where publicly available);
  \item Ecological Impact Assessments (EcIA) for developments within 2\,km of the study area.
\end{itemize}

\dataplaceholder{N} relevant reports were identified, providing ground-truth habitat classification and/or condition data for \dataplaceholder{N PARCELS} parcels within the study area.

\begin{longtable}{p{15mm}p{30mm}p{20mm}p{15mm}p{45mm}}
\caption{Professional ecological survey reports used for validation.}
\label{tab:ecol-reports} \\
\toprule
\textbf{Ref} & \textbf{Title} & \textbf{Author} & \textbf{Date} & \textbf{Relevance} \\
\midrule
\endfirsthead
\midrule
\endhead
\bottomrule
\endlastfoot

ES-01 & \dataplaceholder{TITLE} & \dataplaceholder{AUTH} & \dataplaceholder{DATE} & Phase~1 habitat survey covering \dataplaceholder{N} parcels in northern study area \\
ES-02 & \dataplaceholder{TITLE} & \dataplaceholder{AUTH} & \dataplaceholder{DATE} & Extended Phase~1 and NVC survey for \dataplaceholder{N} parcels \\
ES-03 & \dataplaceholder{TITLE} & \dataplaceholder{AUTH} & \dataplaceholder{DATE} & Preliminary Ecological Appraisal for development site adjacent to study area \\
\end{longtable}

\subsection{Classification Agreement}

\begin{table}[H]
\centering
\caption{Habitat classification agreement: \BDE{} vs professional ecologist surveys.}
\label{tab:class-agreement-ecol}
\begin{tabular}{lrr}
\toprule
\textbf{Metric} & \textbf{\Tzero{} Epoch} & \textbf{\Tone{} Epoch} \\
\midrule
Comparison parcels      & \dataplaceholder{N} & \dataplaceholder{N} \\
Exact class match       & \dataplaceholder{N} (\dataplaceholder{PCT}\%) & \dataplaceholder{N} (\dataplaceholder{PCT}\%) \\
Match at Level 2        & \dataplaceholder{N} (\dataplaceholder{PCT}\%) & \dataplaceholder{N} (\dataplaceholder{PCT}\%) \\
Cohen's Kappa           & \dataplaceholder{K} & \dataplaceholder{K} \\
\bottomrule
\end{tabular}
\end{table}

\begin{confidencebox}{HIGH}
The Level~3 classification agreement of \dataplaceholder{PCT}\% exceeds the 75\% threshold typically considered acceptable for remote sensing-based habitat classification. At the broader Level~2 grouping, agreement exceeds \dataplaceholder{PCT}\%.
\end{confidencebox}

\subsection{Condition Agreement}

\begin{table}[H]
\centering
\caption{Condition assessment agreement: \BDE{} SCI proxy vs field survey condition scores.}
\label{tab:cond-agreement-ecol}
\begin{tabular}{lrr}
\toprule
\textbf{Metric} & \textbf{\Tzero{} Epoch} & \textbf{\Tone{} Epoch} \\
\midrule
Comparison parcels         & \dataplaceholder{N} & \dataplaceholder{N} \\
Exact condition match      & \dataplaceholder{N} (\dataplaceholder{PCT}\%) & \dataplaceholder{N} (\dataplaceholder{PCT}\%) \\
Within $\pm$1 category     & \dataplaceholder{N} (\dataplaceholder{PCT}\%) & \dataplaceholder{N} (\dataplaceholder{PCT}\%) \\
Systematic bias            & \dataplaceholder{DIRECTION} & \dataplaceholder{DIRECTION} \\
Mean condition difference  & \dataplaceholder{DIFF} & \dataplaceholder{DIFF} \\
\bottomrule
\end{tabular}
\end{table}

\begin{confidencebox}{MEDIUM}
The SCI proxy tends to \dataplaceholder{OVER/UNDER}-estimate condition for \dataplaceholder{HABITAT TYPE} habitats by approximately \dataplaceholder{N} category on average. This systematic bias is accounted for in the Monte Carlo uncertainty analysis (Chapter~\ref{ch:delta}).
\end{confidencebox}


\section{University of Derby Thesis Data}

\subsection{Source Description}

\dataplaceholder{AUTHOR} (\dataplaceholder{YEAR}) conducted a habitat condition survey of the Ashleyhay area as part of an undergraduate dissertation at the University of Derby, School of Biological and Environmental Sciences. The thesis, titled ``\dataplaceholder{THESIS TITLE}'', provides:

\begin{itemize}[nosep]
  \item NVC quadrat data for \dataplaceholder{N} grassland parcels within the study area;
  \item Condition assessments using the Natural England Common Standards Monitoring (CSM) methodology;
  \item Species lists for \dataplaceholder{N} hedgerow segments;
  \item Photographic evidence of habitat condition at the survey date.
\end{itemize}

\begin{confidencebox}{MEDIUM}
The thesis was conducted in \dataplaceholder{YEAR}, which falls between the two assessment epochs. The thesis data therefore provides a temporal midpoint validation rather than a direct comparison at either epoch. Habitat condition trajectories implied by the thesis data are compared against the \BDE{}'s interpolated trajectory.
\end{confidencebox}

\subsection{NVC Correspondence}

The thesis provides NVC community assignments for surveyed quadrats. These are compared against the \BDE{}'s UKHab classifications, using the published NVC-to-UKHab correspondence table \citep{ukhab_v2}:

\begin{table}[H]
\centering
\caption{NVC community correspondence: thesis data vs \BDE{} UKHab classification.}
\label{tab:nvc-correspondence}
\begin{tabular}{p{25mm}p{30mm}p{30mm}p{25mm}}
\toprule
\textbf{Thesis NVC} & \textbf{Expected UKHab} & \textbf{BDE UKHab} & \textbf{Agreement} \\
\midrule
MG5 (Centaureo--Cynosuretum)  & g3a5 & \dataplaceholder{CLASS} & \dataplaceholder{Y/N} \\
MG6 (Lolium--Cynosurus)       & g4    & \dataplaceholder{CLASS} & \dataplaceholder{Y/N} \\
U1 (Festuca--Agrostis--Rumex) & g1a   & \dataplaceholder{CLASS} & \dataplaceholder{Y/N} \\
W10 (Quercus--Pteridium--Rubus) & w1f & \dataplaceholder{CLASS} & \dataplaceholder{Y/N} \\
\dataplaceholder{NVC} & \dataplaceholder{UKHAB} & \dataplaceholder{CLASS} & \dataplaceholder{Y/N} \\
\bottomrule
\end{tabular}
\end{table}

\subsection{Condition Trajectory Comparison}

Where the thesis provides condition data for parcels also assessed by the \BDE{} at both epochs, the implied condition trajectory (change between the thesis date and each epoch) can be examined:

\begin{figure}[H]
\centering
\begin{tikzpicture}
  \draw[dreamlab-light,fill=dreamlab-light,rounded corners] (0,0) rectangle (12,7);
  \node[dreamlab-grey,font=\sffamily\large] at (6,3.5) {Figure placeholder: Condition trajectory comparison};
  \node[dreamlab-grey,font=\sffamily\footnotesize] at (6,2.5) {../output/figures/condition\_trajectory\_comparison.png};
\end{tikzpicture}
\caption{Condition trajectory comparison for parcels with thesis ground-truth data. Each line represents one parcel; dots indicate \BDE{} assessments at \Tzero{} and \Tone; triangles indicate thesis assessment at \dataplaceholder{YEAR}. Consistent trajectories are shown in green; inconsistent trajectories in red.}
\label{fig:condition-trajectory}
\end{figure}


\section{Priority Habitat Inventory Cross-Check}

\subsection{PHI Coverage}

The Natural England Priority Habitat Inventory v2.3 identifies \dataplaceholder{N} priority habitat polygons within the study area, covering \dataplaceholder{AREA}\,ha (\dataplaceholder{PCT}\% of the study area). These polygons carry habitat type and, in some cases, condition information that can be compared against the \BDE{} classification.

\subsection{PHI vs BDE Classification}

\begin{table}[H]
\centering
\caption{Classification agreement between PHI designations and \BDE{} assignments.}
\label{tab:phi-agreement}
\begin{tabular}{lrr}
\toprule
\textbf{Metric} & \textbf{Value} & \textbf{Notes} \\
\midrule
PHI polygons in study area       & \dataplaceholder{N} & --- \\
Matched to BDE parcels           & \dataplaceholder{N} & One-to-many matching where PHI polygon spans multiple parcels \\
Habitat type agreement           & \dataplaceholder{PCT}\% & At UKHab Level~2 \\
False positives (BDE, not PHI)   & \dataplaceholder{N} & Parcels classified as priority habitat by BDE but not in PHI \\
False negatives (PHI, not BDE)   & \dataplaceholder{N} & PHI polygons not classified as priority habitat by BDE \\
\bottomrule
\end{tabular}
\end{table}

\begin{confidencebox}{MEDIUM}
The PHI is known to be incomplete, particularly for small or degraded habitat patches. False positives (BDE identifying priority habitat not in PHI) may represent genuine priority habitat not yet captured in the inventory. False negatives warrant investigation as potential classifier errors.
\end{confidencebox}


\section{Systematic Bias Assessment}

Across all three validation sources, the following systematic patterns emerge:

\begin{table}[H]
\centering
\caption{Summary of systematic biases identified through cross-examination.}
\label{tab:bias-summary}
\begin{tabular}{p{40mm}p{25mm}p{60mm}}
\toprule
\textbf{Bias} & \textbf{Magnitude} & \textbf{Impact on Delta} \\
\midrule
Condition over-estimation in woodland habitats & $+$\dataplaceholder{N} category (mean) & Inflates woodland BU at both epochs; effect on delta is partially cancelled \\
Confusion between modified and neutral grassland & \dataplaceholder{PCT}\% of grassland parcels & May over-state grassland improvement delta \\
Hedgerow species-richness classification uncertainty & \dataplaceholder{PCT}\% of segments & Affects hedgerow unit distinctiveness scores \\
\bottomrule
\end{tabular}
\end{table}


\section{Cross-Examination Conclusions}

\begin{confidencebox}{HIGH}
The cross-examination confirms that the \BDE{} pipeline produces habitat classifications and condition assessments within acceptable tolerances of independent ground-truth data. The identified systematic biases are of small magnitude and are partially self-cancelling in the delta calculation. All biases are incorporated into the Monte Carlo uncertainty model, ensuring that the reported confidence intervals adequately capture the validation discrepancies.
\end{confidencebox}

Key findings:
\begin{enumerate}[nosep]
  \item Habitat classification overall agreement exceeds \dataplaceholder{PCT}\% at Level~3 and \dataplaceholder{PCT}\% at Level~2;
  \item Condition assessment agrees within $\pm$1 category for \dataplaceholder{PCT}\% of validated parcels;
  \item NVC community correspondence from the thesis data confirms the plausibility of the dominant habitat types identified by the classifier;
  \item The magnitude of the reported delta is robust to the identified biases.
\end{enumerate}
